% !TEX root = main.tex
\subsection{Thrust 3: Evidence-Guided Trust Containment and Recovery.}
\label{sb:t3}

\textbf{Problem Statement.}
Most of an embedded device's security-critical and functional logic runs not in the Secure world but as a single, monolithic \emph{Normal-world} RTOS application -- protocol stacks, drivers, RTOS scheduling, and control logic sharing one privilege level and one flat address space. Compartmentalization reduces the blast radius of a compromise here by restricting each component to only the code, data, and peripherals it requires~\cite{clements2018aces,kim2018securing,khan2023ec}, but on resource-constrained MCUs this least-privilege goal is at odds with the hardware available to enforce it: the MPU supports only a small, fixed number of regions, forcing coarse compartments or costly per-transition reconfiguration, and that reconfiguration can itself violate real-time deadlines. Beyond this hardware mismatch, spatial isolation and control-flow protection are rarely composed safely, and correctness under nested, preempting interrupts and DMA-capable peripherals is largely unanalyzed. Even where isolation and control-flow protection both hold, the data that legitimately crosses a compartment boundary remains unconstrained: existing systems restrict \emph{how} compartments may communicate but not \emph{what} they may do with the state they exchange, leaving a confused-deputy gap that a compromised compartment can exploit without violating control-flow policy.
Finally, when a violation is detected, existing systems respond only locally and expose no evidence to a remote verifier, so a contained fault is remotely indistinguishable from a propagated compromise.

This thrust therefore asks the following fundamental question: \textbf{\textit{How can a resource-constrained MCU's Normal-world RTOS application be compartmentalized so that both control flow and the data crossing compartment boundaries are authenticated under real-time and interrupt constraints, and so that a violation produces evidence a remote verifier can use to contain and selectively recover the affected compartment?}} Addressing this question requires new compiler-driven partitioning algorithms that jointly reason about MPU region constraints, real-time budgets, and DMA-capable peripherals; authenticated data-flow policies; interrupt-safe transition mechanisms whose correctness can be formally verified; and evidence-guided recovery protocols that connect to the runtime attestation developed in Thrust~2.

\textbf{Related Work.}
Prior MCU compartmentalization systems partition Normal-world RTOS applications via program-dependence-graph analysis, but none jointly addresses real-time cost, interrupt safety, and data-flow integrity. ACES~\cite{clements2018aces} and EC~\cite{khan2023ec} rely on static dependence analysis, but MPU reconfiguration costs tens of percent of runtime and neither analyzes nested-interrupt correctness; EC further keeps compartments and its own monitor in privileged mode rather than dropping to unprivileged execution, trading CPU-level least privilege for lower porting effort~\cite{khan2023ec}. MINION~\cite{kim2018securing} instead uses coarse, whole-thread views to preserve real-time behavior. None constrains data crossing a compartment boundary: ACES retains write access for up to 80\% of store instructions to a shared variable under its most restrictive policy~\cite{clements2018aces}, and EC's shared-type qualifier grants every compartment full read/write access~\cite{khan2023ec}, leaving a confused-deputy gap exploitable without violating control-flow policy~\cite{chen2005non}. D-Box~\cite{mera2022d} narrows this for DMA channels alone, leaving CPU-writable descriptor fields unprotected.
None of these systems exposes evidence for remote containment and recovery~\cite{clements2018aces,kim2018securing,mera2022d,khan2023ec}, a gap Thrust~3 addresses.

\begin{comment}
% --- Original, longer Problem Statement (pre-shortening), kept for reference. ---
\textbf{Problem Statement.}
Embedded system's bulk of a device's firmware wjtb most security-critical logic and functional logic is not hosted on the Secure-world.
The majority of embedded and cyber-physical application code -- protocol stacks, drivers, RTOS scheduling, and control logic -- continues to execute as a single, monolithic \emph{Normal-world} (non-secure) RTOS application that shares one privilege level and one flat address space, and that must be defended in its own right rather than assumed away by the existence of a separate secure domain.
Compartmentalization reduces the blast radius of a compromise in this Normal-world RTOS code by restricting each firmware component to only the code, data, and peripherals it requires, rather than trusting the entire application with the full ambient privilege of the execution environment~\cite{clements2018aces,kim2018securing,khan2023ec}.
On resource-constrained MCUs, however, this least-privilege goal is fundamentally at odds with the hardware available to enforce it.
The Memory Protection Unit (MPU) supports only a small, fixed number of size- and alignment-constrained regions, forcing compilers and runtimes to either merge unrelated code and data into the same region -- weakening isolation -- or to reconfigure the MPU on every cross-compartment transition, which existing systems show can cost tens of percent of runtime even after aggressive optimization~\cite{clements2018aces,khan2023ec}. Real-time and RTOS-based MCUs face an additional constraint: frequent privilege-mode switching to reconfigure protection state can itself violate deadline guarantees, which is why systems such as MINION~\cite{kim2018securing} restrict themselves to coarse, whole-thread memory views rather than fine-grained code compartments in order to preserve real-time behavior.

A second limitation is that spatial isolation and architectural control-flow protection are rarely composed safely. Existing compartmentalization systems either omit control-flow authentication entirely, allowing a compromised compartment to still reuse any gadget reachable within its own code region~\cite{kim2018securing}, or rely on shared runtime/switching code and shared cross-compartment call metadata that enlarge the set of valid branch targets across compartment boundaries~\cite{clements2018aces,khan2023ec}. Because MPU reconfiguration is not atomic with respect to interrupts, an exception or fault can also be delivered while protection state is only partially updated, and even formally verified enforcement mechanisms have had to acknowledge that a compromised RTOS can mask or override the very debug exceptions used to police this transition~\cite{khan2023ec}. None of these systems analyzes correctness under \emph{nested}, preempting interrupts specifically: ACES does not model interrupts at all; MINION extends its memory-view switching to interrupt entry but does not discuss what happens if a higher-priority interrupt preempts a lower-priority one, or preempts the view-switching routine itself; and EC avoids the question for compartment switches by masking ordinary interrupts for the duration of each switch rather than proving switches correct under nesting. Even production RTOS memory protection fares no better here: FreeRTOS-MPU and Zephyr's memory-domain support run interrupt handlers in privileged mode with unrestricted access to all memory, exempting ISRs from protection entirely rather than reasoning about their safe interleaving with compartment transitions. DMA-capable peripherals compound the problem: because DMA controllers operate outside the CPU's MPU-mediated path, most compartmentalization systems ignore DMA-based attacks altogether, and the one system that does address DMA restricts its capability model to a single peripheral class rather than treating all bus masters uniformly~\cite{mera2022d}.

A third limitation is that even where spatial isolation and control-flow authentication are both enforced, the data that legitimately crosses a compartment boundary remains unconstrained and unverified. Existing systems restrict \emph{how} compartments may communicate, but not \emph{what} they may do with the state they exchange. ACES's own evaluation shows that once a global variable is placed in a shared data region, a median of up to 80\% of the store instructions in the application binary retain write access to it even under its most restrictive policy~\cite{clements2018aces}, and EC's \texttt{shared} type qualifier grants a variable full read/write accessibility to every compartment in the system rather than to only the specific compartments and directions an interaction actually requires~\cite{khan2023ec}. Authenticating that a call reached its intended target compartment does not constrain the values carried across that call: a compromised compartment can still supply a well-formed but malicious argument, or overwrite shared state that another, uncompromised compartment later reads, without violating any control-flow policy. For example, an authorization flag staged in a shared region -- such as the PIN-validity flag in ACES's own PinLock case study~\cite{clements2018aces} -- can be set directly by any compartment holding write access to that region, bypassing the credential check it is meant to gate without triggering any control-flow violation; this is precisely the class of non-control-data attack shown to be a realistic threat even when control-flow integrity holds~\cite{chen2005non}. This is also a form of the classical confused-deputy problem, and existing systems largely acknowledge rather than solve it. ACES restricts and validates only the \emph{locations} of compartment transitions, leaving the values that flow through them unchecked~\cite{clements2018aces}; MINION's own discussion identifies shared data between processes as a residual attack surface that a compromised process can exploit to subvert another~\cite{kim2018securing}; and EC explicitly defers both confused-deputy mitigation and data-flow-aware partitioning to future work~\cite{khan2023ec}. D-Box shows that a capability-based access model can narrow this problem for a single class of shared resource -- DMA channels -- by binding each channel to one task rather than granting ambient access~\cite{mera2022d}, but this protection does not extend to the CPU-side shared memory that general-purpose compartments use to communicate, where the confused-deputy gap remains open: a compartment with legitimate permission to queue a transfer can still overwrite another task's in-flight DMA descriptor -- its source, destination, and length fields, which are ordinary CPU-writable shared memory -- to redirect where the transfer writes, an attack outside the scope of D-Box's channel-ownership model. Recent work outside the MCU compartmentalization literature shows that runtime data-flow integrity (DFI) can be made practical for real-time systems by exploiting slack between average execution time and worst-case execution time reservations~\cite{294543}, but this result targets Cortex-A/Linux-capable platforms with hardware (ARM Memory Tagging Extension) and profiling infrastructure unavailable on resource-constrained Cortex-M MCUs, protects a single task's own instrumented code rather than authenticating data flow across separate compartments, and enforces locally without producing evidence a remote verifier can act on.

A fourth and largely unaddressed limitation is that when a violation is detected, existing systems respond only locally -- typically by faulting or halting the offending compartment -- and expose no evidence to a remote verifier about which compartment was compromised, what it attempted to do, or whether the remaining compartments can still be trusted~\cite{clements2018aces,kim2018securing,mera2022d,khan2023ec}. Consequently, an operator has no principled way to distinguish a localized fault that can be contained and recovered from a compromise that has propagated beyond its compartment boundary, and no way to selectively restore and re-establish trust in only the affected component rather than the entire device.

This thrust therefore asks the following fundamental question: \textbf{\textit{How can a resource-constrained MCU's Normal-world, RTOS-managed application be fine-grained and interrupt-safe compartmentalized so that cross-compartment control-flow transitions and the data they carry are authenticated alongside spatial isolation, and so that runtime violations produce evidence that lets a remote verifier localize the violation, contain it to the affected compartment, and selectively recover and re-attest the rest of the system?}} Addressing this question requires new compiler-driven partitioning algorithms that jointly reason about MPU region constraints, real-time budgets, and DMA-capable peripherals; authenticated data-flow policies that constrain what compartments may read and write across a shared boundary; interrupt-safe transition mechanisms whose correctness can be formally verified; and evidence-guided recovery protocols that connect the local isolation enforced here to the runtime attestation developed in Thrust~2.
\end{comment}




\textbf{Task 3.1: Real-Time- and Data-Flow-Aware Compartmentalization with Interrupt-Safe Transitions.}
We will develop a novel compiler-driven partitioning algorithm that assigns application code and data to compartments using a dependence-graph analysis, with two objectives not considered by existing work: MPU region-budget feasibility under a real-time deadline, and directional separation of a shared variable's readers from its writers.
This dependence analysis is the foundation the remaining tasks in this thrust build their enforcement and recovery mechanisms on.
As shown in Figure~\ref{fig:t31}, the partitioner consumes this dependence graph together with each task's worst-case-execution-time (WCET) budget and produces compartments annotated with directional MPU regions.


\begin{wrapfigure}{r}{0.52\textwidth}
	\centering
	%\vspace{-10pt}
	\begin{tikzpicture}[
			font=\tiny,
			box/.style={draw, rounded corners, align=center, minimum height=0.6cm, inner sep=2pt, fill=white},
			state/.style={draw, rounded corners, align=center, minimum height=0.45cm, minimum width=1.6cm, inner sep=1.5pt, fill=white},
			fstate/.style={draw, rounded corners, dashed, align=center, minimum height=0.45cm, minimum width=1.5cm, inner sep=1.5pt, fill=red!6},
			arr/.style={-{Latex[length=1.4mm]}, thick},
			farr/.style={-{Latex[length=1.4mm]}, thick, dashed, red!55!black},
		]
		% ---- Left column: partitioning pipeline ----
		\node[box, fill=green!14, text width=2.3cm] (pdgwcet) {Dependence graph (control-flow, data, peripherals) + WCET budget};
		\node[box, fill=green!24, text width=2.3cm, below=3mm of pdgwcet] (part) {Real-time- \& data-flow-aware partitioner};
		\node[box, fill=green!8, text width=2.3cm, below=3mm of part] (comp) {Compartments $C_1..C_n$: directional MPU regions (R/O~/~R/W)};
		\draw[arr] (pdgwcet) -- (part);
		\draw[arr] (part) -- (comp);

		% ---- Right column: transition monitor FSM, top-aligned with left column ----
		\node[state, fill=green!6, anchor=north west] (idle) at ($(pdgwcet.north east)+(15mm,0)$) {IDLE};
		\node[state, below=4mm of idle] (revoke) {REVOKE $C_i$};
		\node[state, below=6mm of revoke] (grant) {GRANT $C_j$};
		\node[state, fill=green!6, below=4mm of grant] (commit) {COMMIT: run $C_j$};
		\node[fstate, text width=1.5cm, right=6mm of grant] (fault) {FAULT: access $\leq \min(C_i,C_j)$};

		\draw[arr] (idle) -- node[right]{\texttt{SVC}} (revoke);
		\draw[arr] (revoke) -- node[left]{$C_i$ revoked} (grant);
		\draw[arr] (grant) -- node[right]{$C_j$ granted} (commit);
		\draw[farr] ($(revoke)!0.5!(grant)$) -- node[above,sloped]{NMI/HardFault} (fault);
		\node[font=\tiny, below=1mm of commit] (nextnote) {(next \texttt{SVC} $\to$ IDLE)};
		\node[font=\tiny\bfseries, below=1mm of nextnote, text width=3.6cm, align=center] {Transition monitor FSM, e.g.\ $C_i \to C_j$};
	\end{tikzpicture}
	\vspace{-0.3cm}
	\caption{\textit{Task 3.1}: compartments come from a PDG analysis and real-time budgets (left); every transfer, including DMA commits, drives the transition monitor's verified FSM (right), which stays safe by construction under an NMI/HardFault mid-transition.}
	\label{fig:t31}
	\vspace{-10pt}
\end{wrapfigure}

We will make every cross-compartment control transfer pass through a single trusted transition monitor reached only via \texttt{SVC}.
At boot, the monitor is configured to run at the highest priority any compartment can ever request, and that configuration is then locked against modification by any running compartment.
 %only modifiable by the initial trusted boot sequence -- never a running compartment -- can change it afterward.
Compared to EC's~\cite{khan2023ec} approach of masking interrupts at every switch, which enables a compromised RTOS to divert control flow by priority manipulation, our design fixes non-preemptibility at configuration time, making it resistant to such runtime subversion.
A locked highest priority does not, however, shield the monitor from Cortex-M's \texttt{NMI} and \texttt{HardFault} exceptions, which sit above any configurable priority level and cannot be masked.
To this end, we sequence each MPU region update in fail-safe order, revoking the outgoing compartment's permissions before granting the incoming compartment's, so that any protection state an \texttt{NMI} observes mid-transition is never more permissive than either compartment's legitimate access, only potentially less.
We will formally verify this ordering property and the transition monitor's control logic, following the model-checking and proof-assistant approaches.
Because verifying arbitrary nesting depth is a substantially harder target than the interrupt-masked properties, we will first bound the verified depth to the platform's configured \texttt{NVIC} priority levels, extending to unbounded nesting only if the resulting state space remains tractable.
Finally, we will treat DMA transfer requests as cross-domain calls rather than unilateral peripheral writes: a \textit{trusted DMA broker} validates a compartment's requested source, destination, and length against that compartment's own memory boundaries before committing the descriptor to hardware (Figure~\ref{fig:t31}), protecting the CPU-writable descriptor fields ordinary DMA validation leaves exposed.
Because validating every transfer individually would reintroduce the per-transfer cost DMA exists to avoid, we will amortize this cost across repeated, same-shaped transfers by caching a transfer's shape and cheaply re-authorizing subsequent transfers that match it.

\begin{wrapfigure}{r}{0.48\textwidth}
	\centering
	\vspace{-8pt}
	\begin{tikzpicture}[
			font=\tiny,
			box/.style={draw, rounded corners, align=center, minimum height=0.7cm, inner sep=2.5pt, fill=white},
			mech/.style={draw, rounded corners, align=center, minimum height=0.7cm, inner sep=2pt, fill=gray!12},
			dec/.style={draw, diamond, aspect=2.2, align=center, inner sep=0.5pt, fill=blue!6},
			arr/.style={-{Latex[length=1.6mm]}, thick},
		]
		\node[box, fill=blue!10, text width=2.6cm] (edge) {Def-use edge from Task 3.1's dependence graph};
		\node[dec, below=3mm of edge, text width=2.2cm] (amb) {Statically unambiguous?};
		\node[box, fill=blue!16, text width=1.5cm, below left=4mm and -3mm of amb] (static) {Static MPU region: R/O reader, R/W writer};
		\node[mech, text width=1.5cm, below right=4mm and -3mm of amb] (pa) {PA-tagged write: writer ID, def-site, epoch};
		\node[box, fill=blue!8, text width=2.6cm, below=5mm of amb, yshift=-8mm] (check) {Checked at read, or lazily via evidence stream};
		\draw[arr] (edge) -- (amb);
		\draw[arr] (amb) -- node[above,sloped,pos=0.55]{yes} (static);
		\draw[arr] (amb) -- node[above,sloped,pos=0.55]{no} (pa);
		\draw[arr] (pa) -- (check);
	\end{tikzpicture}
	\caption{\textit{Task 3.2}: unambiguous dependencies get a static, directional MPU region for free; ambiguous ones fall back to a PA-tagged write, checked at read or via the evidence stream.}
	\label{fig:t32}
	\vspace{-10pt}
\end{wrapfigure}
\textbf{Task 3.2: Ambiguity-Guided Data-Flow Integrity.}
Task~3.2 targets the data that legitimately crosses a compartment boundary, authenticating which compartment may write a given shared location rather than leaving all shared state uniformly accessible to every compartment with region access.
The central challenge is enforcing this on Cortex-M, which has no hardware memory-tagging support, without paying full-instrumentation cost for locations the compiler can already resolve statically.

We address this challenge by letting Task~3.1's dependence graph decide, for each shared location, whether it can be protected with a static MPU region or needs a runtime check.
Where the graph proves a single compartment is the only possible writer, we enforce that with a directional MPU region: one compartment gets read-write access, every other compartment gets read-only (or no) access, and the hardware enforces this boundary at no added cost.
Where the graph cannot narrow the writer down to one compartment, due to pointer aliasing or interrupt-interleaved updates (Figure~\ref{fig:t32}), we fall back to a lightweight signature built on ARM's PA hardware: each write computes an authenticator over the writer's identity, the write location, and a monotonically increasing epoch, using the \texttt{PACGA} instruction and the same PA keys Thrust~2 uses for control-flow evidence.
A read recomputes this authenticator and compares it against the small set of candidate-writer keys the dependence graph could not rule out; a compromised compartment outside that set cannot forge a matching tag, and the epoch prevents replaying an old one.
This bounds the confused-deputy gap by constraining which compartment may write a given ambiguous location, while spending its cost only where ambiguity cannot be statically resolved, so overhead scales with program ambiguity rather than with the number of shared variables or the volume of execution.
\begin{comment}
% --- Original Task 3.2 text (pre-rewrite), kept for reference. ---
Building on the same ambiguity-driven philosophy we use for control-flow evidence in Task~2.1, we will treat each data dependency Task~3.1's partitioner exposes across a compartment boundary as either statically resolvable or ambiguous, and enforce it accordingly rather than applying one uniform mechanism to all shared state.
As illustrated in Figure~\ref{fig:t32}, where the dependence analysis can prove a shared location has a single reachable writer for a given reader, we enforce that relationship for free using directional (read-only / read-write) MPU permissions, automating -- and pushing below task granularity -- a distinction that production RTOS memory-protection APIs such as FreeRTOS-MPU and Zephyr's memory domains already expose architecturally but require a developer to configure manually, one task at a time.

Where the analysis cannot statically bound the reaching writer -- due to pointer aliasing or interrupt-interleaved updates, the same conditions that force ACES's own over-approximation~\cite{clements2018aces} -- we fall back to a dynamic mechanism only for that specific dependency, and only among the small candidate set of compartments the analysis could not rule out, rather than trusting any code with region access, as ACES, EC, and MINION do today. Because Cortex-M has no hardware memory-tagging support, each such location is paired with a compact shadow-metadata slot: on write, the compartment computes an authenticator over its own compartment identity, the specific def-site, and the current transition epoch using the \texttt{PACGA} generic-authentication instruction and the same Pointer Authentication key infrastructure Thrust~2 already uses for control-flow measurement, and stores it in that slot. A read -- or, to avoid per-access overhead, the evidence report -- recomputes and compares the authenticator against each key in the candidate-writer set; because producing a valid tag requires the corresponding compartment's key, which untrusted code cannot access, a compromised compartment that can overwrite the metadata slot but does not hold that key still cannot forge a tag that verifies. We will enforce the epoch as a monotonically increasing counter checked for freshness on every read, so that a compromised compartment cannot replay a stale-but-genuinely-valid tag captured from an earlier, legitimate write. This bounds, without fully eliminating, the confused-deputy gap identified in the Problem Statement -- it constrains which compartment may successfully write a given ambiguous location, not whether the value that compartment legitimately chooses to write is itself correct -- and it does so without either the blanket instrumentation cost of classical software DFI~\cite{castro2006securing} or the Cortex-A-only hardware assumptions of opportunistic real-time DFI~\cite{294543}. Because this mechanism spends its cost only where the dependence analysis is genuinely uncertain, we expect its overhead to scale with program ambiguity rather than with the number of shared variables or the volume of execution, mirroring the design goal already validated for control-flow evidence in \enola.
\end{comment}

\textbf{Task 3.3: Evidence-Guided Trust Containment and Recovery.}
We will unify the three enforcement points developed above: MPU faults, transition-monitor violations, and data-flow tag mismatches into a single compact evidence report that extends Task~2.1's encoding with a compartment identifier, the faulting program counter, and the last authenticated transition.
On detecting a violation, the runtime will revoke only the faulting compartment's MPU regions and invalidate its entry in the secure task-to-key mapping.
The same evidence report lets the verifier confirm that no \emph{other} compartment accepted a transition or data-flow tag from the faulting compartment after the point of compromise.
This guarantee is scoped to what Tasks~3.1 and~3.2 authenticate, it does not claim the faulting compartment caused no harm before detection .


\begin{comment}
% --- Figure 9 (fig:t33), removed: largely restates the Task 3.3 prose (violation ->
% localized evidence -> contain/verify -> selective recovery) without adding structure
% not already stated as complete sentences in the text. Kept here for reference. ---
\begin{wrapfigure}{r}{0.48\textwidth}
	\centering
	%\vspace{-8pt}
	\begin{tikzpicture}[
			font=\tiny,
			box/.style={draw, rounded corners, align=center, minimum height=0.7cm, inner sep=2.5pt, fill=white},
			mech/.style={draw, rounded corners, align=center, minimum height=0.7cm, inner sep=2pt, fill=gray!12},
			arr/.style={-{Latex[length=1.6mm]}, thick},
		]
		\node[box, fill=violet!10, text width=2.6cm] (viol) {Violation: MPU fault, FSM violation, or PA tag mismatch};
		\node[mech, text width=2.6cm, below=3mm of viol] (evid) {Localized evidence: compartment ID, fault PC, last transition};
		\node[box, fill=violet!16, text width=1.5cm, below left=3mm and 5mm of evid] (contain) {Contain: revoke MPU regions \& keys};
		\node[box, fill=violet!16, text width=1.5cm, below right=2mm and 5mm of evid] (verify) {Verify: no accepted evidence elsewhere};
		\node[box, fill=violet!8, text width=2.6cm, below=4mm of $(contain)!0.5!(verify)$] (recover) {Selective re-attest \& recover via Thrust~1};
		\draw[arr] (viol) -- (evid);
		\draw[arr] (evid) -- (contain);
		\draw[arr] (evid) -- (verify);
		\draw[arr] (contain) -- (recover);
		\draw[arr] (verify) -- (recover);
	\end{tikzpicture}
	\caption{\textit{Task 3.3}: violations from any enforcement point in Tasks 3.1--3.2 are unified into one evidence report driving localized containment and selective, Thrust-1-based re-attestation of only the affected compartment.}
	\label{fig:t33}
	\vspace{-10pt}
\end{wrapfigure}
\end{comment}

We will then recover the affected compartment rather than the whole device: using Thrust~1's domain-construction and attestation mechanisms, we rebuild compartment's trust measurement, re-verify, and re-admit it into the running system, avoiding a full reboot.
The central challenge is minimizing state revocation lists and per-compartment key storage across compartments to prevent replay attacks or reuse a revoked key, while keeping this bookkeeping within real-time budgets.
